\section{Related Work}\label{sec:related}

\textbf{Trajectory datasets.} NGSIM provides camera-derived trajectories at 0.1-s intervals~\cite{fhwa2016ngsim}. highD records naturalistic highway traffic from drones, including vehicle dimensions and classes~\cite{krajewski2018highd}. I-24 MOTION extends observation over a 4.2-mi freeway corridor using roadside cameras~\cite{gloudemans2023motion}. VTT complements them with onboard motion and control signals, platform-linked size and powertrain labels, and explicit tests on unseen car models. Its source logs come from the open-source driving-automation ecosystem used by comma2k19~\cite{schafer2018comma2k19} and ADAS-TO~\cite{wang2026adasto}. Public onboard datasets also support external checks: VED records vehicle energy and driving data~\cite{oh2019ved}, while OpenACC captures commercial adaptive-cruise-control behaviour~\cite{makridis2021openacc}.

\textbf{Vehicle attributes from motion.} GPS-based classifiers distinguish passenger cars and trucks using engineered features~\cite{sun2013gps} or convolutional networks~\cite{dabiri2020gps}. Naturalistic and instrumented studies report differences between electric and combustion vehicles in driver adaptation~\cite{helmbrecht2014ev}, adaptive-cruise-control operation~\cite{geffen2026ev} and regenerative-braking phases~\cite{liu2025paai}. We examine whether these differences support powertrain identification, alongside finer footprint-based size classes, on unseen car models.

\textbf{Identity cues and generalisation.} CAN signals can identify drivers from complete drives~\cite{enev2016fingerprinting} or individual turns~\cite{hallac2016turn}. If windows from the same driver or car model appear in training and testing, such cues can inflate type-recognition scores. This is a form of shortcut learning~\cite{geirhos2020shortcut}. VTT separates evaluation on unseen models from evaluation on unseen drivers of known models, and estimates uncertainty by resampling car models.
